\documentclass[a4paper]{article}
\usepackage{amsfonts}
\usepackage{amsmath}
\usepackage{amssymb}
\usepackage{graphicx}   

\newcommand{\Bgcot}{\operatorname{Bgcot}}

\begin{document}
  
\noindent \large Auteur: Elaine de Bie \\
\noindent \large Studierichting: Informatica\\
\noindent \large Oefeningenreeks 2B nr. 8.4\\


\medskip


\normalsize


$\Bgcot 2 = \Bgcot 3 + \Bgcot 7$\\

$\Leftrightarrow \cot(\Bgcot 2) = \cot(\Bgcot 3 + \Bgcot 7)$\\

Stel dan $ \alpha = \Bgcot 3 $ en $\beta = \Bgcot 7$\\

$\Leftrightarrow \cot(\Bgcot 2) = \cot(\alpha + \beta)$\\

Gebruik de cotangent optelformule: 
\[\cot(\alpha + \beta) = \dfrac{\cot\alpha \cot\beta -1}{\cot\beta + \cot\alpha}\]

$\Leftrightarrow \cot(\Bgcot 2) = \dfrac{\cot\alpha \cot\beta -1}{\cot\beta + \cot\alpha}$\\

Invullen van $\cot\alpha = 3$ en $\cot\beta = 7$:\\

$\Leftrightarrow \cot(\Bgcot 2) = \dfrac{3 \cdot 7 -1}{7 + 3}$\\

$\Leftrightarrow \cot(\Bgcot 2) = \dfrac{21 - 1}{7 + 3}$\\

$\Leftrightarrow \cot(\Bgcot 2) = \dfrac{20}{10}$\\

$\Leftrightarrow \cot(\Bgcot 2) = 2$\\

$\Leftrightarrow 2 = 2$\\




\begin{thebibliography}{999}
%\bibitem{a} Referentie
\end{thebibliography}
\end{document}
